\begin{foldable} % Appendix overview
    In this extended text, we provide our implementation details (\ref{sec:app01-training}), a complementary study on alternative optimization objectives for our image priors (\ref{sec:app02-experiments-multiobjective}), and an extension to an iterative framework (\ref{sec:app03-extension}). Explicit results to the figures in the main text are also provided for completeness and reproducibility (\ref{sec:app04-explicit}).
\end{foldable}



\section{Training Details} \label{sec:app01-training}
\renewcommand{\thefigure}{A.\arabic{figure}}
\renewcommand{\thetable}{A.\arabic{table}}
\renewcommand{\theequation}{A.\arabic{equation}}
\setcounter{figure}{0}
\setcounter{table}{0}
\setcounter{equation}{0}

\subsection{Datasets} \label{ssec:app01-training-datasets}
\begin{foldable} % Datasets details
    In our experiments (Sec.~\ref{sec:04-experiments}), we evaluate our method on eight benchmark image datasets: MNIST \citep{lecun2002gradient}, USPS \citep{hull1994database}, SVHN \citep{netzer2011reading}, FMNIST \citep{xiao2017fashion}, CIFAR10, CIFAR100 \citep{krizhevsky2009learning}, TinyImageNet \citep{le2015tiny}, Imagenette \citep{howard2020imagenette}.
    All datasets are public and easily downloadable with standard libraries.
    In terms of image content, MNIST, USPS, and SVHN contain small images of digits.
    FMNIST is also a small-sized dataset of fashion items.
    The images in CIFAR10, CIFAR100, TinyImageNet, and Imagenette all contain diverse scenes and objects, e.g., animals, plants, food, vehicles, and household goods, to name a few.
    Notably, TinyImageNet and Imagenette have higher resolution: 64$\times$64 for TinyImageNet, and approximately one magnitude larger than that for Imagenette (ranging from 27$\times$80 to 4368$\times$2912 --- which we reported the average dimension).
    Of all the benchmarks, only CIFAR100 and TinyImageNet have abundant classes (100 and 200, respectively), while the others have 10 classes.
    Regarding the dataset size, USPS and Imagenette are relatively smaller.
    We summarize their details in Table~\ref{tab:appendix-training-datasets}.
\end{foldable}

\begin{table}[ht] % tab:appendix-training-datasets
    \caption{Details of the image benchmark datasets.}
    \centering
    \begin{tabular}{c|c|c|cc|c|c}
        \hline                         & \textbf{Dataset} & \textbf{Dimension} & \multicolumn{2}{c|}{\textbf{Size}} & \textbf{No.} & \textbf{Content}  \\
        & & & train & test & \textbf{classes} & \\
        \hline
        \multirow{4}{*}{Simple} & MNIST  & 1$\times$28$\times$28 & \phantom{0}60 K & 10 K           & \phantom{0}10 & Digits \\
                                & USPS   & 1$\times$16$\times$16 & \phantom{00}9 K & \phantom{0}2 K & \phantom{0}10 & Digits \\
                                & SVHN   & 3$\times$32$\times$32 & \phantom{0}73 K & 26 K           & \phantom{0}10 & Digits \\
                                & FMNIST & 1$\times$28$\times$28 & \phantom{0}60 K & 10 K           & \phantom{0}10 & Fashion items \\
        \hline
        \multirow{4}{*}{Complex} & CIFAR10      &                    3$\times$32$\times$32 & \phantom{0}50 K &           10 K & \phantom{0}10 & Diverse \\
                                 & CIFAR100     &                    3$\times$32$\times$32 & \phantom{0}50 K &           10 K &           100 & Diverse \\
                                 & TinyImageNet &                    3$\times$64$\times$64 &           100 K &           10 K &           200 & Diverse \\
                                 & Imagenette   & 3$\times$408$\times$472$^\blacktriangle$ & \phantom{00}9 K & \phantom{0}4 K & \phantom{0}10 & Diverse \\
        \hline
        \multicolumn{5}{l}{$^\blacktriangle$: The average dimension of Imagenette is reported.} \\
    \end{tabular}
    \label{tab:appendix-training-datasets}
\end{table}


\subsection{Teacher and students training} \label{ssec:app01-teacherstudents}
\subsubsection{Augmentation} \label{sssec:app01-teacherstudents-augment}
\begin{foldable} % Augmentation
    To augment the training data (real images for the teacher, synthetic images for students), we design a basic augmentation pipeline.
    Firstly, if the images are very small (in MNIST, USPS, FMNIST), we resize them to at least $32\times32$.
    Then the images undergo random shifting by 1/8 of their dimension, and random horizontal flipping unless they are from a digits dataset (MNIST, USPS, SVHN).
    Finally, we randomly jitter the color in the images for more image variety.
    % For the high-resolution Imagenette, we efficiently perform $4\times$ downscale at storage and $4\times$ upscale at runtime.
\end{foldable}

\subsubsection{Optimization and scheduling} \label{sssec:app01-teacherstudents-optimization}
\begin{foldable} % Optimization and scheduling
    To train the Teacher for baseline, and the Students in our method, we use a stochastic gradient descent optimizer with momentum of 0.9, weight decay of $5\times10^{-4}$, and batch size of 100.
    On datasets with small black-and-white images (MNIST, USPS, FMNIST), we use a fixed learning rate of $10^{-2}$ and train for 100 epochs.
    On other datasets, we set the learning rate to $10^{-1}$ at the beginning and gradually decrease by two orders of magnitude with a scheduler to enable optimal convergence for 200 training epochs.
\end{foldable}


\subsection{Optimization of Image Priors} \label{ssec:app01-imagepriors}
\subsubsection{Diverging filter} \label{sssec:app01-teacherstudents-optimization-divgfilter}
\begin{foldable} % Diverging filter principles
    We provide more details on the formulation of the diverging filter $f_{\text{divg}}$ (Eq.~\ref{eq:synthesis-diverging-filter}) proposed in the Synthesis phase (Sec.~\ref{ssec:03-framework-synthesis}).
    To fulfill its role in contrasting opposite-type regions in a mask $m\in[0, 1]^{C\times H\times W}$, we expect $f_{\text{divg}}$ to satisfy the following properties:
    \begin{enumerate}
        \item We start by picking a brightness cutoff threshold $\beta$ to separate the negative ($x<\beta$) and postive regions ($x\ge\beta$) for any pixel $x\in m$.
            The extrema and the threshold value must be fixed after filtering, \ie, $f_{\text{divg}}\left(\cdot; \beta\right)$ must go through $\left\{ (0,0), (\beta,\beta), (1,1)\right\}, \forall\beta$.
        \item To preserve the order of intensity across the whole mask, it is necessary that $f_{\text{divg}}$ is an increasing function, \ie, $f_{\text{divg}}^{\prime}\left(x;\beta\right)>0,\forall x,\beta\in\left[0,1\right]$.
        \item To avoid creating abrupt changes in intensity, \ie, $f_{\text{divg}}$ must be continuously differentiable on $[0,1], \forall\beta$.
        \item To ensure the diverging property emerges at the threshold value $x=\beta$, we require that $f_{\text{divg}}$ has the sharpest increase there, \ie, $f_{\text{divg}}^{\prime}(\cdot;\beta)$ has a maximum at $x=\beta\Leftrightarrow f^{\prime\prime}(x=\beta;\beta)=0,f^{\prime}(x=\beta;\beta)>0$.
    \end{enumerate}
\end{foldable}

\begin{figure}[h] % fig:diverging-filter
    \centering
    \includegraphics[width=0.55\textwidth]{./figures/diverging-filter.pdf}
    \caption{
        Our diverging filter as $y=f_{\text{divg}}\left(x;\beta\right)|_{x\in[0, 1]}$ in Eq.~\ref{eq:synthesis-diverging-filter} with different choices of $\beta\in\left[0,1\right]$.
        Scattered points $\left\{\left(\beta,\beta\right)\right\}$ along the diagonal $y=x$ are also inflection points.
    }
    \label{fig:diverging-filter}
\end{figure}

\begin{foldable} % Diverging filter design
    Accordingly, we design a piecewise function consisting of two continuously differentiable half-parabolas connected at the inflection point $\left(\beta,\beta\right)$.
    One of them has a minimum at the $x=0$ boundary, and vice versa, the other half-parabola has a maximum at the opposite boundary $x=1$, \ie, $f_{\text{divg}}^{\prime}(0;\beta)=f_{\text{divg}}^{\prime}(1;\beta)=0 \forall \beta\in[0, 1]$.
    Solving the two parabolas with the proposed set of constraints indeed gives the solution  we reported $f_{\text{divg}}$ in Eq.~\ref{eq:synthesis-diverging-filter}.
    We plot $f_{\text{divg}}$ with a range of $\beta \in [0, 1]$ values in Fig.~\ref{fig:diverging-filter} to demonstrate its diverging property.
    We believe that this filter design has unexplored potentials in image augmentation for robust learning of vision tasks.
\end{foldable}



\subsubsection{Loading unique embeddings during Contrast phase} \label{sssec:app01-teacherstudents-optimization-uniqueembeddings}
\begin{foldable} % Problem: exploding gradients
    In the Contrast phase, we apply the contrastive method proposed in \citet{fang2021contrastive} to make our data diverse.
    However, we occasionally observed numerical inbstability and synthetic image corruption during contrastive optimization.
    After investigation, we found that the issue stems from how embeddings are loaded for contrastive learning.
    Originally, the embeddings of past inputs $x$ and their corresponding positive views $x^{+}$ are cached to and recalled from an index-free memory bank.
    (The same occurred for negative samples, but it is not linked to the error.)
    As a result, for small datasets, an embedding loaded from the memory bank might be near-duplicated with one extracted at runtime as, by chance, they may originate from the same image and differ by some trivial updates.
    Optimization-wise, contrasting an image against itself is counter-intuitive, which explains the exploding gradients and corruption of synthetic images.
\end{foldable}

\begin{foldable} % Alternative fix
    Therefore, we modify how images are loaded for contrastive optimization.
    Instead of using an index-free memory bank, we store the embeddings of each synthetic images in our dataset by its unique index.
    At contrastive learning runtime, we only load embeddings of images \textbf{not} in the current batch, successfully avoid duplication.
    This modification fixes the numerical instability and synthetic image corruption on our ends.
\end{foldable}


\subsubsection{Instance-discriminator} \label{sssec:app01-teacherstudents-optimization-instancedisc}
\begin{foldable} % Instance-discriminator
    Following \citet{fang2021contrastive}, we use a fully connected instance-discriminator $R$ comprising of two libear layers.
    For each input image $x$, $R$ projects the image embedding $h=S_0^{\mathcal{H}}(x)$ from the primer student backbone $S_0^{\mathcal{H}}$, through one hidden layer, to its own embedding space of 128 dimensions.
    To update $R$ during the Contrast phase (Sec.~\ref{ssec:03-framework-contrast}), we let the contrastive loss $\mathcal{L}_{\text{CT}}$ in Eq.~\ref{eq:loss-contrastive} backpropagate to it parameters.
    It is trained with an Adam optimizer \citep{kingma2014adam} with a learning rate of $10^{-3}$, and its ``betas'' coefficients of (0.5, 0.999).
    As $R$ is shallow, this does not incur a significant computational overhead.
\end{foldable}

\subsubsection{Resources} \label{sssec:app01-teacherstudents-optimization-resources}
\begin{foldable} % Resources
    We perform each standard experiment on a single NVIDIA V100 GPU with 40GB VRAM, 4--16 CPUs, and 4--128 GB RAM, taking around 6--96 hours.
    The required resources depend on the dataset complexity, model size, and synthetic image resolution.
    In an increasing order, their training time rank as: simple datasets (negligible difference), CIFAR10, CIFAR100, Imagenette, and TinyImageNet.
\end{foldable}



\section{Alternative image optimization objectives} \label{sec:app02-experiments-multiobjective}
\renewcommand{\thefigure}{B.\arabic{figure}}
\renewcommand{\thetable}{B.\arabic{table}}
\renewcommand{\theequation}{B.\arabic{equation}}
\setcounter{figure}{0}
\setcounter{table}{0}
\setcounter{equation}{0}


\begin{foldable} % In-class similarity: One-hot loss
    In the following section, we will examine student accuracy when adopting different optimization objectives for image prior optimization, and justify our choice of the contrastive approach for diversity. 
    During the Contrast phase, we employ the contrastive loss $\mathcal{L}_{\text{CT}}$ in Eq.~\ref{eq:loss-contrastive} to make synthetic images more diverse.
    Unlike the diversity objective in our framework, previous BBDFKD methods such as \citet{zhang2022ideal, yuan2024data} attempted to promote in-class similarity and class-balance.
    Specifically, in-class similarity is optimized via the one-hot loss:
    \begin{equation} \label{loss-onehot}
        \mathcal{L}_{\text{OH}}=\mathcal{L_{\text{CE}}}\left(\hat{y}_{S-\text{Soft}},\text{argmax}(\hat{y}_{S-\text{Soft}})\right),
    \end{equation}
    which guides data generation to strengthen the effect of the most prominent class in each sample.
    The one-hot loss is popular in data-free KD \citep{chen2019data}, but as stated in \citet{yuan2024data}, its adversarial nature might cause ``overfitting of synthetic data to the clone model'', resulting in mode collapse and insufficient diversity.
\end{foldable}

\begin{foldable} % Class-balance: Information entropy loss
    The other objective, class-balance, is to ensure the samples are evenly distributed between all classes, enforced by the information entropy loss:
    \begin{equation} \label{loss-infoentropy}
        \mathcal{L}_{\text{IE}}=\sum_{k=0}^{K-1}p_{k}\log p_{k},\text{ with }p_{k}=\mathbb{E}_{x\sim \mathcal{D}}\left[\hat{y}_{S-\text{Soft}}^{\left(k\right)}\right],
    \end{equation}
    where $k$ denotes the $k$-th class probability.
    In our framework, as we have actively balanced the dataset during sampling, class-balance has already been satisfied.
\end{foldable}

\begin{figure}[ht] % fig:abla-losses
    \centering
    \begin{minipage}{0.35\columnwidth}
        \centering
        \includegraphics[width=\columnwidth]{./figures/experiments-abla-losses-mnist-lenet5.pdf}
        {\footnotesize \textbf{(a)} MNIST}
    \end{minipage}
    \hspace{0.05\textwidth}
    \begin{minipage}{0.35\columnwidth}
        \centering
        \includegraphics[width=\columnwidth]{./figures/experiments-abla-losses-cifar10-resnet18.pdf}
        {\footnotesize \textbf{(b)} CIFAR10}
    \end{minipage}
    \caption{Contribution to distillation accuracy (\%) of one-hot loss ($\mathcal{L}_{\text{OH}}$), information entropy loss ($\mathcal{L}_{\text{IE}}$), and contrastive loss ($\mathcal{L}_{\text{CT}}$), on MNIST and CIFAR10.}
    \label{fig:abla-losses}
\end{figure}

\begin{foldable} % Objectives discussion
    To measure how these objectives might affect KD performance, we ablate them using different losses to optimize the image priors dataset $\mathcal{D} = \{x | x \sim \mathcal{X}_\mathcal{P}\}$, on our standard MNIST and CIFAR10 experiments as the base cases.
    In details, we ran all eight binary combinations of using/not using $\mathcal{L}_{\text{OH}}$, $\mathcal{L}_{\text{IE}}$, $\mathcal{L}_{\text{CT}}$ to optimize $\mathcal{D}$, and report in Fig.~\ref{fig:abla-losses}.
    The contribution of each objective is derived by the mean difference in student accuracy between using/not using the corresponding loss, while standard error is aggregated on a per-combination basis.
    We observe $\mathcal{L}_{\text{IE}}$ only has modest contribution, as class-balance has already been achieved during Synthesis.
    Surprisedly, in-class similarity $\mathcal{L}_{\text{OH}}$ is actually adversarially destructive, harming KD performance by a significant amount of $-$2.65\%.
    Meanwhile, contrastive optimization using $\mathcal{L}_{\text{CT}}$ fruitfully helped gain KD performance, notably by $+$0.12\% on MNIST and $+$3.14\%.
    Thus, we focus solely on diversity and employ only the contrastive loss $\mathcal{L}_{\text{CT}}$ to optimize image priors.
\end{foldable}

\begin{foldable} % Outro
    On a side note, \citet{zhang2022ideal, yuan2024data} employed $\mathcal{L}_{\text{OH}}$ and $\mathcal{L}_{\text{IE}}$ on the weights of a generator network, while we employ $\mathcal{L}_{\text{CT}}$ on the parameterized image priors $\mathcal{X}_{\mathcal{P}}$ and an instance-discriminator $R$.
    We speculate that this difference also affects the comparision and requires further investigation.
\end{foldable}



\section{Extension to Iterative Framework} \label{sec:app03-extension}
\renewcommand{\thefigure}{C.\arabic{figure}}
\renewcommand{\thetable}{C.\arabic{table}}
\renewcommand{\theequation}{C.\arabic{equation}}
\setcounter{figure}{0}
\setcounter{table}{0}
\setcounter{equation}{0}

\begin{foldable}
    An interesting extension of our framework is performing \textit{iterative Contrast-Distillation} to further boost student performance.
    Recall that in the Contrast phase (Sec.~\ref{ssec:03-framework-contrast}), we employ an primer student $S_0$ to optimize the dataset $\mathcal{D}$, and train a new student $S$ in the next Distillation phase.
    It naturally emerges that the stronger $S$ possesses even more robust embeddings and can be \textit{recycled for Contrast phase again}.
    Iteratively, we may re-employ the latest student $S_g$ at the $g$\nobreakdash-th generation in the role of a new primer student to optimize $\mathcal{D}$ in the next Contrast phase, and query soft labels from it to train the new student $S_{g+1}$ in the next Distillation phase.
    The objectives for $S_{g+1}$, extended from $\mathcal{L}_S$ (Eq.~\ref{eq:loss-kd-total}), is to match the teacher's label $T(x)$ with $\mathcal{L}_\text{Hard-KD}$, and to match previous student's logits $S_{g}^{\mathcal{Z}}(x)$ with $\mathcal{L}_\text{Soft-KD}$, on samples $x\sim\mathcal{D}$:
    \begin{equation} \label{eq:loss-kd-extension}
        \mathcal{L}_{S_g}
        = \mathcal{L}_{\text{Hard-KD}}+\mathcal{L}_{\text{Soft-KD}}
        = \mathbb{E}_{x\sim\mathcal{D}}\Big[\mathcal{L}_{\text{CE}}\left(\hat{y}_{S_{g+1}-\text{Soft}},\hat{y}_{T-\text{Hard}}\right) + \mathcal{L}_{\text{L1}}\left(S_{g+1}^{\mathcal{Z}}\left(x\right),S_{g}^{\mathcal{Z}}\left(x\right)\right)\Big].
    \end{equation}
\end{foldable}

\begin{figure}[h] % fig:abla-ngenerations
    \centering
    \includegraphics[width=0.75\textwidth]{./figures/experiments-abla-ngenerations.pdf}
    \caption{Distillation accuracy (\%) by generation $g$.}
    \label{fig:abla-ngenerations}
\end{figure}

\begin{foldable}
    From our standard experiments (Sec.~\ref{ssec:04-experiments-standard}, corresponds to $g = 1$), we iterate for more generations to reach $g = 4$ on simple datasets and $g = 2$ on complex datasets due to resource constraints.
    As reported in Fig.~\ref{fig:abla-ngenerations}, even though KD performance converges quickly on simple dataset, it iteratively improves until the last generations with trivial fluctuations.
    These results conclude that our iterative Contrast-Distillation strategy can robustly deliver increasingly better students with more generations, and that generalization is not hurt with longer training.
    Nonetheless, computational cost increases \textit{linearly} but student accuracy does not scale as fast to justify the trade-off.
    We look forward to making this extension even more efficient in future works.
\end{foldable}



\section{Explicit Results} \label{sec:app04-explicit}
\renewcommand{\thefigure}{D.\arabic{figure}}
\renewcommand{\thetable}{D.\arabic{table}}
\renewcommand{\theequation}{D.\arabic{equation}}
\setcounter{figure}{0}
\setcounter{table}{0}
\setcounter{equation}{0}


In this section, we provide explicit numerical results for the figures presented in the main text and above sections of the Appendix for completeness and reproducibility.


\subsection{Increasing synthetic dataset size} \label{ssec:app04-explicit-nsamples}
\begin{foldable}
    We provide explicit student accuracy and standard error for KD by number of samples, which has been visualized and discussed in Fig.~\ref{fig:abla-nsamples}, Sec.~\ref{sssec:04-experiments-ablation-nsamples}.
    The results are grouped by the simple datasets in Table~\ref{tab:abla-nsamples-simple} and complex datasets in Table~\ref{tab:abla-nsamples-complex}.
\end{foldable}

\begin{table}[h] % tab:abla-compression
    \caption{Distillation accuracy (\%) by number of sample $N$ on simple datasets.}
    \centering
    \begin{tabular}{c|c|c|c|c}
        \hline
            $N$
            & \makecell{\textbf{MNIST}   \\ LeNet5}
            & \makecell{\textbf{USPS}    \\ LeNet5}
            & \makecell{\textbf{SVHN}    \\ AlexNet}
            & \makecell{\textbf{FMNIST}  \\ LeNet5} \\
        \hline
        \phantom{0}10 K &          98.20$_{\pm .06}$ &          93.22$_{\pm .43}$ &          94.64$_{\pm .13}$ &                    76.14$_{\pm 2.48}$ \\
        \phantom{0}20 K &          98.59$_{\pm .08}$ &          93.99$_{\pm .24}$ &          95.46$_{\pm .07}$ &                    79.49$_{\pm 1.38}$ \\
        \phantom{0}30 K &          98.68$_{\pm .03}$ &          94.05$_{\pm .08}$ &          95.71$_{\pm .05}$ &          82.67$_{\pm \phantom{0}.45}$ \\
        \phantom{0}50 K &          98.75$_{\pm .03}$ &          94.20$_{\pm .21}$ &          95.94$_{\pm .05}$ &          83.98$_{\pm \phantom{0}.65}$ \\
        \phantom{0}70 K &          98.81$_{\pm .03}$ &          94.25$_{\pm .06}$ &          95.97$_{\pm .02}$ &          84.84$_{\pm \phantom{0}.69}$ \\
                  100 K &          98.78$_{\pm .05}$ &          94.41$_{\pm .08}$ &          96.06$_{\pm .03}$ & \textbf{85.63}$_{\pm \phantom{0}.39}$ \\
                  128 K & \textbf{98.90}$_{\pm .02}$ & \textbf{94.44}$_{\pm .09}$ & \textbf{96.13}$_{\pm .01}$ &          85.62$_{\pm \phantom{0}.06}$ \\
        \hline
        \multicolumn{5}{l}{Best results are in \textbf{bold}.} \\
    \end{tabular}
    \label{tab:abla-nsamples-simple}
\end{table}

\begin{table}[ht] % tab:abla-compression
    \caption{Distillation accuracy (\%) by number of samples $N$ on complex datasets.}
    \centering
    \begin{tabular}{c|c|c|c|c}
        \hline
            $N$ & \textbf{CIFAR10} & \textbf{CIFAR100} & \textbf{TinyImageNet} & \textbf{Imagenette} \\
        \hline
        \phantom{0}10 K &                    63.54$_{\pm 2.64}$ &                    25.85$_{\pm 1.31}$ &          07.88$_{\pm .49}$ &                    50.82$_{\pm 1.65}$ \\
        \phantom{0}20 K &          72.98$_{\pm \phantom{0}.50}$ &          37.48$_{\pm \phantom{0}.48}$ &          12.13$_{\pm .17}$ &                    57.43$_{\pm 1.51}$ \\
        \phantom{0}30 K &          78.18$_{\pm \phantom{0}.60}$ &          42.82$_{\pm \phantom{0}.56}$ &          13.96$_{\pm .42}$ &                    58.98$_{\pm 1.19}$ \\
        \phantom{0}50 K &          83.41$_{\pm \phantom{0}.46}$ &          48.45$_{\pm \phantom{0}.26}$ &          16.49$_{\pm .37}$ &          61.84$_{\pm \phantom{0}.77}$ \\
        \phantom{0}70 K &          85.20$_{\pm \phantom{0}.35}$ &          50.75$_{\pm \phantom{0}.18}$ &          18.69$_{\pm .63}$ &          63.08$_{\pm \phantom{0}.96}$ \\
                  100 K &          86.60$_{\pm \phantom{0}.29}$ &          52.85$_{\pm \phantom{0}.30}$ &          19.99$_{\pm .36}$ &          65.98$_{\pm \phantom{0}.78}$ \\
                  128 K &          87.56$_{\pm \phantom{0}.11}$ &          53.69$_{\pm \phantom{0}.26}$ &          21.07$_{\pm .50}$ & \textbf{66.60}$_{\pm \phantom{0}.72}$ \\
                  256 K & \textbf{89.24}$_{\pm \phantom{0}.15}$ & \textbf{56.84}$_{\pm \phantom{0}.54}$ &          24.46$_{\pm .26}$ &                                     - \\
                  512 K &                                     - &                                     - & \textbf{28.06}$_{\pm .22}$ &                                     - \\
        \hline
        \multicolumn{5}{l}{Best results are in \textbf{bold}.} \\
    \end{tabular}
    \label{tab:abla-nsamples-complex}
\end{table}

\subsection{Students of compressed capacity} \label{ssec:app04-explicit-compression}
We provide in Table~\ref{tab:abla-compression} explicit student accuracy and standard error for KD by compression ratio, which has been visualized and discussed in Fig.~\ref{fig:abla-compression}, Sec.~\ref{sssec:04-experiments-ablation-compress}.

\begin{table}[ht] % tab:abla-compression
    \caption{Distillation accuracy (\%) by compression ratio (CR).}
    \centering
    \begin{tabular}{c|c|c|c|c|c}
        \hline
            {\centering \textbf{CR}}
            & \makecell{\textbf{MNIST}   \\ LeNet5}
            & \makecell{\textbf{USPS}    \\ LeNet5}
            & \makecell{\textbf{SVHN}    \\ AlexNet}
            & \makecell{\textbf{FMNIST}  \\ LeNet5}
            & \makecell{\textbf{CIFAR10} \\ ResNet18} \\
        \hline
                  100\% & 98.59$_{\pm \phantom{0}.08}$ & 94.20$_{\pm \phantom{0}.21}$ & 95.94$_{\pm \phantom{0}.05}$ & 83.98$_{\pm \phantom{0}.65}$ & 83.41$_{\pm \phantom{0}.46}$ \\
        \phantom{0}25\% &           96.79$_{\pm 1.39}$ & 92.42$_{\pm \phantom{0}.08}$ & 95.41$_{\pm \phantom{0}.04}$ & 74.90$_{\pm \phantom{0}.37}$ & 79.23$_{\pm \phantom{0}.51}$ \\
        \phantom{00}4\% &           85.83$_{\pm 1.86}$ & 90.70$_{\pm \phantom{0}.33}$ & 92.51$_{\pm \phantom{0}.07}$ &           52.55$_{\pm 3.79}$ & 64.48$_{\pm \phantom{0}.81}$ \\
        \phantom{00}1\% &           47.72$_{\pm 7.44}$ &           52.43$_{\pm 1.47}$ &           72.00$_{\pm 2.89}$ &           36.85$_{\pm 8.66}$ &           47.39$_{\pm 3.30}$ \\
        \hline
    \end{tabular}
    \label{tab:abla-compression}
\end{table}
        

\subsection{Extension to Iterative Framework} \label{ssec:app03-extension-ngenerations}
\begin{foldable}
    We provide explicit student accuracy and standard error for KD by number of generations, which has been visualized and discussed in Fig.~\ref{fig:abla-ngenerations}, Sec.~\ref{ssec:app03-extension-ngenerations}.
    We ran the experiments for all datasets, but limits to $g$ = 2 on complex dataset due to resource constraints, as shown in Table~\ref{tab:abla-ngenerations}.
    We additionally remark that the student can get iteratively better until the last generations, with some trivial fluctuations.
\end{foldable}

\begin{table}[h] % tab:abla-generation
    \caption{Distillation accuracy (\%) by generation $g$.}
    \centering
    \begin{tabular}{c|c|c|c|c|c|c|c|c}
        \hline
            $g$
            & \makecell{{MNIST}   \\ LeNet5}
            & \makecell{{USPS}    \\ LeNet5}
            & \makecell{{SVHN}    \\ AlexNet}
            & \makecell{{FMNIST}  \\ LeNet5}
            & \makecell{{CIFAR10} \\ ResNet18}
            & \makecell{{CIFAR100} \\ ResNet18}
            & \makecell{{Tiny$^{(*)}$} \\ ResNet18}
            & \makecell{{Nette$^{(**)}$} \\ ResNet18} \\
        \hline
        0 &          98.29$_{\pm .32}$ &          94.04$_{\pm .52}$ &          95.84$_{\pm .10}$ &                    82.38$_{\pm 1.09}$ &          78.17$_{\pm .64}$ &          48.36$_{\pm .19}$ &          22.66$_{\pm .36}$ &          57.67$_{\pm .74}$ \\
        1 &          98.59$_{\pm .08}$ &          94.20$_{\pm .21}$ &          95.94$_{\pm .05}$ &          83.98$_{\pm \phantom{0}.65}$ &          83.41$_{\pm .46}$ &          52.85$_{\pm .30}$ &          28.03$_{\pm .22}$ &          61.84$_{\pm .77}$ \\
        2 &          98.62$_{\pm .12}$ &          94.32$_{\pm .12}$ &          95.95$_{\pm .03}$ &          84.43$_{\pm \phantom{0}.66}$ &          84.50$_{\pm .16}$ & \textbf{53.97}$_{\pm .62}$ & \textbf{30.06}$_{\pm .22}$ & \textbf{62.78}$_{\pm .91}$ \\
        3 &          98.61$_{\pm .08}$ &          94.42$_{\pm .25}$ & \textbf{95.97}$_{\pm .06}$ &          84.50$_{\pm \phantom{0}.46}$ &          85.18$_{\pm .42}$ & - & - & - \\
        4 & \textbf{98.65}$_{\pm .04}$ & \textbf{94.54}$_{\pm .25}$ &          95.94$_{\pm .06}$ & \textbf{84.62}$_{\pm \phantom{0}.86}$ & \textbf{85.23}$_{\pm .39}$ & - & - & - \\
        % 5 &          98.68$_{\pm .02}$ &            94.30$_{\pm .15}$ & \textbf{96.00}$_{\pm .07}$ &          84.44$_{\pm \phantom{0}.98}$ &          85.17$_{\pm .56}$ & - & - & - \\
        % 6 &          98.63$_{\pm .05}$ &            94.40$_{\pm .19}$ &          95.96$_{\pm .06}$ &          84.76$_{\pm \phantom{0}.52}$ & \textbf{85.91}$_{\pm .65}$ & - & - & - \\
        % 7 & \textbf{98.70}$_{\pm .06}$ &            94.60$_{\pm .10}$ &          95.92$_{\pm .07}$ & \textbf{85.28}$_{\pm \phantom{0}.32}$ &          85.77$_{\pm .61}$ & - & - & - \\
        \hline
        \multicolumn{9}{l}{Best results are in \textbf{bold}. $^{(*)}$TinyImageNet, $^{(**)}$Imagenette} \\
    \end{tabular}
    \label{tab:abla-ngenerations}
\end{table}
